\subsection{Notation and conventions}

As a slight abuse of language, we identify a subset $S \subseteq \mathbb{Z}^d$, with the subgraph of $\mathbb{Z}^d$ induced by $S$ consisting of all edges between points in $S$.
In this way, we may speak about graph-theoretic properties, such as degree or connectedness, for a general subset without confusion.
Two subsets $S_1, S_2 \subseteq \mathbb{Z}^d$, are said to be \emph{ambiently isomorphic}, which we denote by $S_1 \cong S_2$, if there exists a graph isomorphism $\phi \colon \mathbb{Z}^d \to \mathbb{Z}^d$ such that $\phi(S_1) = S_2$; this notion is extended to domino configurations in the same way.
